% source: RKF/README.md
\hypertarget{recognition-kernel-framework}{%
\chapter{Recognition Kernel Framework}\label{recognition-kernel-framework}}

\href{https://doi.org/10.5281/zenodo.21760119}{\includesvg{https://zenodo.org/badge/DOI/10.5281/zenodo.21760119.svg}}

\textbf{Author:} Monty Dabas\\
\textbf{ORCID:} 0009-0005-6948-209X

\textbf{Mathematical proof transport, lawful-transition accounting, open-residue classification, and reproducible certificates}

The Recognition Kernel Framework is a technical and mathematical architecture for evaluating whether a recognized structure survives a transformation after lawful transport and declared memory have been accounted for. The remaining open residue is classified and recorded in a reproducible certificate.

The operational engine is the \textbf{Recognition--Null Kernel Engine (RNKE)}. A domain cartridge is represented by

\[
(S_d,R_d,E_d,\tau_d,K_d),
\]

where \texttt{S\_d} is a domain state, \texttt{R\_d} is a recognition operator, \texttt{E\_d} is an evolution or transition, \texttt{tau\_d} is an index or order parameter, and \texttt{K\_d} is a kernel evaluator.

The common execution pattern is:

\begin{verbatim}
state
-> recognition target
-> lawful transition
-> open-residue audit
-> classification
-> technical action
-> certificate
\end{verbatim}

\hypertarget{archival-record}{%
\section{Archival record}\label{archival-record}}

Reviewer-oriented release: \textbf{Version 0.1.0-review}\\
Zenodo DOI: \textbf{10.5281/zenodo.21760119}

Cite the exact release or commit and the certificate package actually reviewed.

\hypertarget{first-public-verification-cartridge}{%
\section{First public verification cartridge}\label{first-public-verification-cartridge}}

This repository begins with a mathematical proof-transport cartridge imported from the pinned RH-Framework source commit:

\begin{verbatim}
c588ded973a395b5fce37c82f670616160979a2e
\end{verbatim}

The imported certificate campaign covers:

\begin{verbatim}
oriented cut scalar
-> quarter-turn structure
-> native Euler exponential and flow
-> native logarithm and powers
-> finite natural arithmetic
-> prime factorization identities
-> Möbius and von Mangoldt identities
-> Dirichlet zeta and Euler-product obligations on Re(s) > 1
\end{verbatim}

Current archived status:

\begin{longtable}[]{@{}
  >{\raggedright\arraybackslash}p{(\columnwidth - 2\tabcolsep) * \real{0.5000}}
  >{\raggedright\arraybackslash}p{(\columnwidth - 2\tabcolsep) * \real{0.5000}}@{}}
\toprule\noalign{}
\begin{minipage}[b]{\linewidth}\raggedright
Package
\end{minipage} & \begin{minipage}[b]{\linewidth}\raggedright
Status
\end{minipage} \\
\midrule\noalign{}
\endhead
\bottomrule\noalign{}
\endlastfoot
\texttt{F00\_F00E\_EULER\_V0\_1} & \texttt{PASS\_F00\_F00E\_RIGOROUS\_COMPUTATIONAL\_AUDIT} \\
\texttt{F00GHI\_LOG\_ARITHMETIC\_ZETA\_V0\_1} & \texttt{PASS\_F00GHI\_LOG\_ARITHMETIC\_ZETA\_AUDIT} \\
Implemented foundation campaign & \texttt{PASS\_IMPLEMENTED\_FOUNDATION\_CERTIFICATE\_CAMPAIGN} \\
\end{longtable}

Canonical campaign SHA-256:

\begin{verbatim}
8728ab0319afe5fb6e6329deaee479cfa93f50bb4422ee863aaaf4f911dd9582
\end{verbatim}

\hypertarget{what-the-pass-means}{%
\section{What the PASS means}\label{what-the-pass-means}}

A PASS means that every declared computational obligation in the implemented package completed successfully, required source pins matched, exact or outward bounds satisfied the package rules, and adversarial negative controls were detected.

A PASS does \textbf{not} certify claims outside the package boundary.

\hypertarget{scientific-boundary}{%
\section{Scientific boundary}\label{scientific-boundary}}

This release does not certify:

\begin{verbatim}
Fourier–Poisson–Gamma–xi completion;
explicit formula or completed-Weil positivity;
active-band source positivity;
eta-zero endpoint positivity;
the Riemann Hypothesis.
\end{verbatim}

Those items remain \texttt{NOT\ CERTIFIED} or \texttt{OPEN} in this software release.

\hypertarget{consolidated-theorem-archive}{%
\section{Consolidated theorem archive}\label{consolidated-theorem-archive}}

The reviewer-facing theorem surface is under \href{theorum/}{\texttt{theorum/}}. It contains the transferred RH-framework capsules, native cut/completion results, the current generator-to-jet capstone chain, and the thermodynamic theorem archive.

The current advanced reading chain is:

\begin{verbatim}
theorum/41_cut_graded_universal_generator_theorem.md
-> theorum/42_cut_graded_lambda_jacobian_tower_theorem.md
-> theorum/43_bilateral_jet_flow_recognition_capstone_theorem.md
-> theorum/44_madhava_smriti_bilateral_jet_flow_closure_theorem.md
-> theorum/thermodynamics/README.md
\end{verbatim}

Use \href{theorum/README.md}{\texttt{theorum/README.md}} for the complete map, \href{theorum/SOURCE_PR_INDEX.md}{\texttt{theorum/SOURCE\_PR\_INDEX.md}} for source ancestry, and \href{theorum/BRANCH_CONSOLIDATION_AUDIT.md}{\texttt{theorum/BRANCH\_CONSOLIDATION\_AUDIT.md}} for the repository-wide branch verification.

The theorem capsules retain their own evidence labels and claim boundaries. Their inclusion in \texttt{main} does not silently upgrade \texttt{LOCAL\ PASS}, \texttt{USER-REPORTED\ PASS}, or \texttt{IMPLEMENTED\ /\ USER\ RUN\ REQUIRED} into a stronger certification state. Apparently even repositories need adults in the room.

\hypertarget{reviewer-route}{%
\section{Reviewer route}\label{reviewer-route}}

\begin{enumerate}
\def\labelenumi{\arabic{enumi}.}
\tightlist
\item
  Read \href{REVIEWER_GUIDE.md}{\texttt{REVIEWER\_GUIDE.md}}.
\item
  Read \href{CLAIM_BOUNDARY.md}{\texttt{CLAIM\_BOUNDARY.md}}.
\item
  Inspect \href{CERTIFICATE_INDEX.md}{\texttt{CERTIFICATE\_INDEX.md}}.
\item
  Read the consolidated theorem map in \href{theorum/README.md}{\texttt{theorum/README.md}}.
\item
  Reproduce the campaign using \href{REPRODUCE.md}{\texttt{REPRODUCE.md}}.
\item
  Compare generated JSON hashes with the archived certificate artifacts.
\item
  Inspect theorem sources, verifier obligations, and negative controls directly.
\end{enumerate}

\hypertarget{terminology}{%
\section{Terminology}\label{terminology}}

The public umbrella name is \textbf{Recognition Kernel Framework}. The operational engine is \textbf{Recognition--Null Kernel Engine (RNKE)}. Recognition-Seam Calculus, EMK, UGD, RMG, and ECL are mathematical, geometric, computational, dynamic, and certificate-runtime layers or representations used within the broader architecture.

\hypertarget{provenance-policy}{%
\section{Provenance policy}\label{provenance-policy}}

Certified source files are preserved byte-for-byte under their historical paths and namespaces. Historical names such as \texttt{rh\_framework} remain inside the imported cartridge because changing them would change source hashes and invalidate the archived certificate identity.

\hypertarget{public-review-scope}{%
\section{Public-review scope}\label{public-review-scope}}

This repository publishes technical material for scientific and engineering review. It certifies only the packages, source pins, obligations, tolerances, and artifacts listed in the certificate index. Public availability does not imply that every research claim in related repositories has been verified.

\begin{flushright}\small\texttt{RKF/README.md}\end{flushright}
